\chapter{Globus Sensation}

\section*{Definition}
Globus sensation is the persistent or intermittent perception of a lump, foreign body, or tightness in the throat in the absence of a true obstructive lesion, typically midline between the thyroid notch and the sternal notch.

\section*{Pathophysiology}
Increased tension of the cricopharyngeal or thyrohyoid muscles, often triggered by gastroesophageal reflux, laryngopharyngeal irritation, or psychological stress, alters sensory feedback from the pharyngeal mucosa. Visceral hypervigilance and central sensitization amplify the perception of normal pharyngeal sensations.

\section*{Causes}
\begin{itemize}
  \item \textbf{Gastroenterologic:} Laryngopharyngeal reflux (LPR), gastroesophageal reflux disease (GERD), esophageal dysmotility, hiatal hernia
  \item \textbf{Psychogenic:} Anxiety disorders, panic disorder, somatization, major depression, grief reaction
  \item \textbf{Otolaryngologic:} Chronic rhinosinusitis with postnasal drip, allergic rhinitis, tonsilloliths, Eagle syndrome (elongated styloid process)
  \item \textbf{Miscellaneous:} Thyroid enlargement (goiter, nodule), cervical osteophytes, medication-induced xerostomia, cricopharyngeal spasm
  \item \textbf{Neoplastic (rare):} Oropharyngeal or esophageal carcinoma (consider if progressive dysphagia, odynophagia, or weight loss present)
\end{itemize}

\section*{When to See a Doctor}
See your doctor if:
\begin{itemize}
  \item Progressive dysphagia or odynophagia (painful swallowing)
  \item Unintentional weight loss, hoarseness for \textgreater{}3 weeks, or hemoptysis
  \item Neck mass, lymphadenopathy, or otalgia (referred ear pain)
  \item Sensation is unilateral or precisely localized to one spot
  \item Onset after age 50 or in the setting of known head/neck cancer history
\end{itemize}

\section*{Related Symptoms}
\begin{itemize}
  \item Dysphagia, throat clearing, chronic cough, hoarseness, heartburn, regurgitation, excessive salivation, xerostomia, anxiety, chest tightness
\end{itemize}
\textbf{Important:} Globus sensation is present in up to 45\% of the general population and is rarely the sole presenting symptom of malignancy, but a thorough history for red-flag symptoms is essential before making the diagnosis of exclusion.

\section*{Self-Care / Home Management}
\begin{itemize}
  \item Drink warm liquids or swallow saliva deliberately (avoid repetitive dry swallowing)
  \item Elevate the head of the bed by 6--8 inches for suspected reflux
  \item Avoid throat clearing (use a sip of water instead); practice relaxed breathing
  \item Eliminate trigger foods: caffeine, spicy foods, citrus, chocolate, alcohol, carbonated beverages
  \item Use saline nasal rinses if postnasal drip is present
\end{itemize}

\section*{How Your Doctor Will Evaluate You}
\begin{itemize}
  \item Flexible nasolaryngoscopy to visualize the hypopharynx and larynx for inflammation, masses, or pooling of secretions
  \item Barium swallow (esophagram) or videofluoroscopic swallow study if structural lesions or dysmotility suspected
  \item Empiric trial of high-dose proton-pump inhibitor (PPI) for 4--8 weeks if LPR is suspected
  \item Esophageal manometry or pH monitoring if symptoms persist despite PPI trial
  \item Neck ultrasound if thyroid enlargement or palpable abnormality
\end{itemize}

\section*{Treatment Options}
\begin{itemize}
  \item \textbf{Reflux-related:} High-dose PPI (e.g., omeprazole 20--40 mg BID) for 8--12 weeks, dietary modification, and head-of-bed elevation
  \item \textbf{Psychogenic:} Cognitive behavioral therapy (CBT), selective serotonin reuptake inhibitors (SSRIs) if anxiety or depression is present
  \item \textbf{Speech--language therapy:} Laryngeal relaxation exercises, diaphragmatic breathing, and biofeedback for cricopharyngeal tension
  \item \textbf{Reassurance:} Explain the benign nature of the condition and advise against repetitive self-examination or throat clearing, which perpetuates the symptom
\end{itemize}

\clearpage
